% P11 paper module: R22 positive strong-gauge angle defect and fixed-vector criterion.

\subsection{Positive strong-gauge angle defect and fixed-vector criterion}
\label{subsec:o3u-angle-defect}

Keep the odd-sector notation at fixed $0<R<S<T_0<U$ from
Subsections~\ref{subsec:o3n-gauge-intertwining},
\ref{subsec:o3s-relative-polar}, and~\ref{subsec:o3t-crosspolar}.  Thus
\[
 W:=W_{R,S,-}^{[T_0]},
 \qquad
 V_U:=U_SWU_R^*,
 \qquad
 \Gamma_U:=W^*V_U,
\]
\[
 \mathfrak P_U:=U_SW-WU_R,
 \qquad
 Z_U:=U_S^*WU_R,
 \qquad
 \Omega_U:=Z_U^*W=U_R^*\Gamma_UU_R.
\]
R15 identifies $\Gamma_U\to I$ strongly with $V_U\to W$ strongly.  The purpose of
this subsection is to encode that strong gate by a positive uniformly bounded
fixed-vector observable.  The moving-gauge operator built from $Z_U-W$ is retained,
but is not promoted to the R15 strong gate.

\begin{proposition}[Positive angle defect for the R15 strong gauge]
\label{prop:o3u-true-angle}
Define
\[
 \boxed{
 \mathscr G_U:=(V_U-W)^*(V_U-W).
 }
\tag{O3U.1}\label{eq:o3u-G-def}
\]
Then
\[
 \boxed{
 \mathscr G_U=2I-\Gamma_U-\Gamma_U^*.
 }
\tag{O3U.2}\label{eq:o3u-G-gamma}
\]
Moreover
\[
 \boxed{0\le\mathscr G_U\le4I,}
\tag{O3U.3}\label{eq:o3u-G-bound}
\]
and for every fixed source vector $f$,
\[
 \boxed{
 \langle\mathscr G_Uf,f\rangle
 =\|(V_U-W)f\|^2.
 }
\tag{O3U.4}\label{eq:o3u-G-rayleigh}
\]
\end{proposition}

\begin{proof}
Since $V_U$ and $W$ are isometries,
\[
\begin{aligned}
 (V_U-W)^*(V_U-W)
 &=V_U^*V_U-V_U^*W-W^*V_U+W^*W\\
 &=2I-\Gamma_U^*-\Gamma_U,
\end{aligned}
\]
which proves~\eqref{eq:o3u-G-gamma}.  Positivity is immediate from
\eqref{eq:o3u-G-def}.  Also $\|V_U-W\|\le2$, hence
$\|\mathscr G_U\|\le4$ and therefore~\eqref{eq:o3u-G-bound}.  Equation
\eqref{eq:o3u-G-rayleigh} follows from the definition.
\end{proof}

\begin{proposition}[Dense-core criterion for strong gauge compatibility]
\label{prop:o3u-dense-core}
Let $\mathcal C$ be any fixed dense subspace of the $R$-source Hilbert space.  Then
\[
 \boxed{
 \Gamma_U\xrightarrow[s]{}I
 \quad\Longleftrightarrow\quad
 \langle\mathscr G_Uf,f\rangle\longrightarrow0
 \quad\text{for every }f\in\mathcal C.
 }
\tag{O3U.5}\label{eq:o3u-dense-core}
\]
In particular, failure of the R15 strong gauge condition is certified by one fixed
vector and one subsequence whenever
\[
 \boxed{
 \limsup_{n\to\infty}
 \langle\mathscr G_{U_n}f,f\rangle>0
 }
\tag{O3U.6}\label{eq:o3u-fixed-witness}
\]
for some fixed $f$ and $U_n\to\infty$.
\end{proposition}

\begin{proof}
By Proposition~\ref{prop:o3n-gauge-criterion},
\[
 \Gamma_U\xrightarrow[s]{}I
 \quad\Longleftrightarrow\quad
 V_U\xrightarrow[s]{}W.
\]
The forward implication in~\eqref{eq:o3u-dense-core} follows from
\eqref{eq:o3u-G-rayleigh}.  Conversely suppose the quadratic forms tend to zero on
$\mathcal C$.  Then
\[
 \|(V_U-W)f\|\to0
 \qquad(f\in\mathcal C).
\]
For arbitrary $x$ and $f\in\mathcal C$,
\[
 \|(V_U-W)x\|
 \le2\|x-f\|+\|(V_U-W)f\|,
\]
because $\|V_U-W\|\le2$ uniformly.  Taking $U\to\infty$ and then using density of
$\mathcal C$ proves strong convergence on the whole source space.  Statement
\eqref{eq:o3u-fixed-witness} is the contrapositive fixed-vector obstruction.
\end{proof}

\begin{remark}[Smooth odd source core]
\label{rem:o3u-smooth-core}
The previously established TC0 graph-core reduction supplies
$C_c^\infty((-R,R))_{\rm odd}$ as a dense core of the odd P11 source space.  Hence the
criterion~\eqref{eq:o3u-dense-core} may be tested on fixed smooth odd source vectors;
no $U$-dependent maximizing vector is needed for a positive strong-convergence proof.
Conversely, an operator-norm lower bound carried only by $U$-dependent directions does
not by itself disprove strong convergence.
\end{remark}

The cross-polar defect of O3t produces a second positive operator.  It is norm-equivalent
to $\mathscr G_U$ but is expressed in a moving source gauge.

\begin{proposition}[Moving-gauge angle defect]
\label{prop:o3u-moving-angle}
Define
\[
 \boxed{
 D_U:=(Z_U-W)^*(Z_U-W).
 }
\tag{O3U.7}\label{eq:o3u-D-def}
\]
Then
\[
 \boxed{
 D_U=\mathfrak P_U^*\mathfrak P_U
 =2I-\Omega_U-\Omega_U^*
 =U_R^*\mathscr G_UU_R.
 }
\tag{O3U.8}\label{eq:o3u-D-identities}
\]
Consequently
\[
 \boxed{
 \|D_U\|
 =\|\mathscr G_U\|
 =\|\mathfrak P_U\|^2
 =\|V_U-W\|^2.
 }
\tag{O3U.9}\label{eq:o3u-D-norm}
\]
If
\[
 \Lambda_U:=(I-WW^*)U_SW,
\]
then the R20 orthogonal decomposition also gives the operator identity
\[
 \boxed{
 D_U
 =\Lambda_U^*\Lambda_U
 +U_R^*(\Gamma_U-I)^*(\Gamma_U-I)U_R.
 }
\tag{O3U.10}\label{eq:o3u-D-positive-split}
\]
\end{proposition}

\begin{proof}
O3t gives $Z_U-W=-U_S^*\mathfrak P_U$, hence
\[
 D_U=\mathfrak P_U^*\mathfrak P_U.
\]
Also
\[
 V_U-W=\mathfrak P_UU_R^*,
\]
so
\[
 \mathscr G_U
 =U_R\mathfrak P_U^*\mathfrak P_UU_R^*,
\]
which proves $D_U=U_R^*\mathscr G_UU_R$.  Substituting
$\Omega_U=U_R^*\Gamma_UU_R$ into~\eqref{eq:o3u-G-gamma} gives
$D_U=2I-\Omega_U-\Omega_U^*$.  Equation~\eqref{eq:o3u-D-norm} follows from unitary
invariance and
$\|T^*T\|=\|T\|^2$.

Finally Proposition~\ref{prop:o3s-relative-decomposition} gives
\[
 \mathfrak P_U
 =\Lambda_U+W(\Gamma_U-I)U_R
\]
with orthogonal ranges.  Taking the product with its adjoint on the source side removes
the cross terms and yields~\eqref{eq:o3u-D-positive-split}.
\end{proof}

\begin{proposition}[Moving-unitary conjugacy does not preserve fixed-vector strong asymptotics]
\label{prop:o3u-moving-countermodel}
There is an abstract gauge system with fixed isometry $W$ for which
\[
 D_n\xrightarrow[s]{}0
\]
while
\[
 \mathscr G_n\not\xrightarrow[s]{}0.
\]
Thus the pointwise strong behavior of $D_U$ cannot be promoted by abstract unitary
conjugacy to the R15 strong gauge condition.
\end{proposition}

\begin{proof}
Let the source and target be $\ell^2(\mathbb N)$ and take $W=I$.  Put
\[
 V:=I-2P_{e_1},
\]
the reflection in the first coordinate.  Then
\[
 \mathscr G_n:=(V-I)^*(V-I)=4P_{e_1}
\]
for every $n$, so $\mathscr G_n$ does not converge strongly to zero.

Choose unitaries $U_{R,n}$ with
\[
 U_{R,n}^*e_1=e_n
\]
and set
\[
 U_{S,n}:=VU_{R,n}.
\]
Then
\[
 V_n=U_{S,n}WU_{R,n}^*=V
\]
for every $n$, while
\[
 \mathfrak P_n
 =U_{S,n}W-WU_{R,n}
 =(V-I)U_{R,n}.
\]
Therefore
\[
 D_n
 =\mathfrak P_n^*\mathfrak P_n
 =U_{R,n}^*(4P_{e_1})U_{R,n}
 =4P_{e_n}
 \xrightarrow[s]{}0.
\]
This proves the claimed non-promotion.  The example is a topology firewall only; it is
not a counterexample to the concrete P11 metric family.
\end{proof}

\begin{remark}[R22-D strong-topology firewall]
\label{rem:o3u-moving-firewall}
At each fixed $U$, $D_U$ and $\mathscr G_U$ are unitarily conjugate and therefore have
the same norm and spectrum.  Since the conjugating unitary $U_R$ depends on $U$, this
pointwise equivalence does not identify their fixed-vector strong limits.  Thus
$D_U$ is a legitimate norm diagnostic and an exact positive form of the O3t cross-polar
defect, but the primary fixed-vector observable for the R15 strong gate is
$\mathscr G_U$.
\end{remark}

The true angle defect also admits a direct metric-product representation.

\begin{proposition}[Metric representation of the strong-gauge angle defect]
\label{prop:o3u-metric-representation}
Using
\[
 U_S=X_SA_S^{-1/2},
 \qquad
 U_R^*=A_R^{-1/2}X_R^*,
\]
one has
\[
 \boxed{
 V_U=X_SA_S^{-1/2}WA_R^{-1/2}X_R^*,
 }
\tag{O3U.11}\label{eq:o3u-V-metric}
\]
and
\[
 \boxed{
 \Gamma_U
 =W^*X_SA_S^{-1/2}WA_R^{-1/2}X_R^*.
 }
\tag{O3U.12}\label{eq:o3u-Gamma-metric}
\]
Consequently
\[
 \boxed{
 \begin{aligned}
 \mathscr G_U
 ={}&2I
 -W^*X_SA_S^{-1/2}WA_R^{-1/2}X_R^*\\
 &-X_RA_R^{-1/2}W^*A_S^{-1/2}X_S^*W.
 \end{aligned}
 }
\tag{O3U.13}\label{eq:o3u-G-metric}
\]
\end{proposition}

\begin{proof}
The first identity is the polar formulas inserted into
$V_U=U_SWU_R^*$.  Left compression by $W^*$ gives~\eqref{eq:o3u-Gamma-metric}; taking
the adjoint and substituting into~\eqref{eq:o3u-G-gamma} gives
\eqref{eq:o3u-G-metric}.
\end{proof}

\begin{remark}[Inverse-square-root information firewall]
\label{rem:o3u-inverse-root-firewall}
Equation~\eqref{eq:o3u-G-metric} does not bypass the R17--R21 information barrier.
Direct control of $\langle\mathscr G_Uf,f\rangle$ for a fixed vector $f$ requires the
relative orientation of the factors $A_R^{-1/2}$ and $A_S^{-1/2}$ inside the cross
product.  Equivalently, evaluating $V_Uf$ requires control of
\[
 A_R^{-1/2}X_R^*f
 \quad\text{and then}\quad
 X_SA_S^{-1/2}W A_R^{-1/2}X_R^*f.
\]
The existing rank-one fixed-pair Gram asymptotics do not determine this inverse-square-root
orientation; Proposition~\ref{prop:o3n-rankone-insufficiency} and the O3t open problem
remain in force.
\end{remark}

\begin{openproblem}[R22-F: fixed-vector P11 strong-gauge asymptotics]
\label{open:o3u-angle-defect}
On the fixed smooth odd source core, determine whether
\[
 \boxed{
 \langle\mathscr G_Uf,f\rangle
 =\|(V_U-W)f\|^2
 \longrightarrow0
 }
\tag{O3U.14}\label{eq:o3u-open-target}
\]
for every fixed $f$, or construct a fixed $f$ and a subsequence $U_n\to\infty$ for
which the limsup is positive.  By Proposition~\ref{prop:o3u-dense-core}, the first
alternative is exactly the R15 strong polar-gauge condition
$\Gamma_U\to I$.

This open problem concerns only the polar gauge component.  Even if it is resolved
positively, the full future transport
\[
 W_{R,S,-}^{[U]}=U_SQU_R^*
\]
also contains the modulus comparison $Q-W$.  No strong terminal-transport conclusion
and no global Object~X construction is asserted here.
\end{openproblem}
